\documentclass{exam-zh}
\usepackage{edu-explanation}

\examsetup{
  question/show-answer = true,
  fillin/show-answer = true,
  solution/show-solution = show-stay,
}

\begin{document}

\eduSectionTitle{黄金分割比：三组比值，三步列式}

\begin{eduProblemBlock}[例题 1]
一条装饰线全长 $30\text{ cm}$，按黄金分割分成较长、较短两段。
求较长一段的长度。

\end{eduProblemBlock}





\begin{eduAuxItem}{思路导航}
第一步：确定已知边和要求的边$\;\to\;$第二步：确定这两条边的比值$\;\to\;$第三步：确定乘还是除\end{eduAuxItem}




\begin{eduSubQuestionItem}{例题 1}
求较长一段的长度。\end{eduSubQuestionItem}

\begin{eduDualExplain}
  \begin{eduExplainLeft}
    \eduColumnTitle{\eduIconBulb}{小贴士}
        \eduSideHint{比值要有方向}{“长比全”必须写成 $\dfrac{L}{T}$，不能只记“长、全都是 $0.618$”。
}
        \eduSideHint{已知分母就乘}{在 $\dfrac{L}{T}=k$ 中，已知的是分母 $T$，所以分子 $L=T\times k$。
}
  \end{eduExplainLeft}

  \begin{eduExplainRight}
    \eduColumnTitle{\eduIconBook}{解答}
      \eduExplainStep{1}{第一步：确定已知边和要求的边}{已知全长 $T=30\text{ cm}$，要求长段 $L$。
}
      \eduExplainStep{2}{第二步：确定这两条边的比值}{记短段为 $S$，长段为 $L$，全长为 $T$。三组两两比值为
\[
  \frac{S}{L}=\frac{\sqrt5-1}{2},\qquad
  \frac{L}{T}=\\ frac{\sqrt5-1}{2},\qquad
  \frac{S}{T}=\frac{3-\sqrt5}{2}.
\]
本题是长段与全长，所以使用
\[
  \frac{L}{T}=\frac{\\ sqrt5-1}{2}.
\]
}
      \eduExplainStep{3}{第三步：确定乘还是除}{已知分母 $T$，所以用全长乘比值：
\[
  L=30\times\frac{\sqrt5-1}{2}
  =15(\sqrt5-1)\text{ cm}.
\]
}
  \end{eduExplainRight}
\end{eduDualExplain}




\begin{eduProblemBlock}[例题 2]
一段窗格按黄金分割分成长、短两段，其中较短一段长 $6\text{ m}$。
求整段窗格的长度。

\end{eduProblemBlock}





\begin{eduAuxItem}{思路导航}
第一步：确定已知边和要求的边$\;\to\;$第二步：确定这两条边的比值$\;\to\;$第三步：确定乘还是除\end{eduAuxItem}




\begin{eduSubQuestionItem}{例题 2}
求整段窗格的长度。\end{eduSubQuestionItem}

\begin{eduDualExplain}
  \begin{eduExplainLeft}
    \eduColumnTitle{\eduIconBulb}{小贴士}
        \eduSideHint{直接用“短比全”}{已知和所求正好是短段、全长，直接使用 $\dfrac{S}{T}=\dfrac{3-\sqrt5}{2}$，不必先求长段。
}
        \eduSideHint{已知分子就除}{在 $\dfrac{S}{T}=q$ 中，已知的是分子 $S$，所以分母 $T=S\div q$。
}
  \end{eduExplainLeft}

  \begin{eduExplainRight}
    \eduColumnTitle{\eduIconBook}{解答}
      \eduExplainStep{1}{第一步：确定已知边和要求的边}{已知短段 $S=6\text{ m}$，要求全长 $T$。
}
      \eduExplainStep{2}{第二步：确定这两条边的比值}{短比全是：
\[
  \frac{S}{T}=\frac{3-\sqrt5}{2}.
\]
}
      \eduExplainStep{3}{第三步：确定乘还是除}{已知分子 $S$，所以用短段除以比值：
\[
  T=6\div\frac{3-\sqrt5}{2}
  =3(3+\sqrt5)\text{ m}.
\]
}
  \end{eduExplainRight}
\end{eduDualExplain}




\begin{eduSummaryBlock}{\eduIconPin}{每题固定三步}
\begin{eduBulletList}
  \item \textbf{第一步：}确定已知边和要求的边是哪两个。  \item \textbf{第二步：}确定这两条边的比值，并写清分子、分母。  \item \textbf{第三步：}已知分母就乘比值；已知分子就除以比值。\end{eduBulletList}
\end{eduSummaryBlock}




\end{document}
